\section{Introduction}

NLIP fixes how users and agents exchange messages~\cite{ecma430,ecma431}. It says nothing
about what an agent does between receiving a request and answering it, and that is where it
reads untrusted text and acts on it. An email assistant reads the user's inbox and has
permission to send mail and delete files. Text in a third-party MCP server's tool
description~\cite{mcp,invariant2025} or in an incoming email~\cite{greshake2023} can tell
the assistant to forward a security code to an outsider. The user never sees that text, and
the action is carried out with their authority.

Prior defenses either constrain what the agent may do, such as a privilege policy or a
separate planning model~\cite{progent,camel}, or check each action against the user's
task~\cite{taskshield,llamafirewall}. We adopt the second idea and fit it to the constraints
of an NLIP security extension (SE). The function sees only the NLIP message, the structured
tool call and its result, its own session state, and the deployer's policy. It never sees the
agent's reasoning.

\paragraph{Contributions.} (1)~\sys, a Type-3 SE at the tool boundary: its NLIP ingress part
records the user's request as the only trusted statement of the task, and its tool-boundary
part judges every side-effecting call against that request and scans every tool result
(\S\ref{sec:design}); it plugs in through the framework's SE configuration and a middleware,
with no attack-specific rules. (2)~A realistic NLIP application: an email, calendar and drive
assistant over MCP with Google OAuth sign-in as an ingress SE (\S\ref{sec:app}). (3)~A paired
end-to-end evaluation on AgentDojo~\cite{agentdojo} test tasks: the same tool-poisoning cases
and benign tasks, run with and without \sys{} on two agent models. It shows whether attacks
that succeed without the guard are stopped with it, and whether benign tasks still complete,
along with false alarms, cost, precision and recall (\S\ref{sec:results}); the scanner is also measured on its own, since AgentDojo's data
injections barely work on these models.
Code, data, a demo and every command are at
\url{https://github.com/MCP-Tool-Guard/ToolGuard} (MIT license).
